\chapter{TÓM LƯỢC}
\indent\indent Trong kỷ nguyên chuyển đổi số, việc tự động phân loại dữ liệu di sản văn hóa phi vật thể (ICH) giữ vai trò quan trọng trong hoạt động bảo tồn và quảng bá văn hóa, song các phương pháp CNN truyền thống chỉ sử dụng dữ liệu hình ảnh vẫn gặp nhiều hạn chế do đặc thù dữ liệu đa phương thức và mức độ tương đồng cao về trang phục, nghi thức và bối cảnh lễ hội. Để giải quyết vấn đề này, tôi đề xuất một kiến trúc học sâu mới kết hợp Graph Attention Network dùng đồ thị đa phương thức (MGAT) với FastKAN (mạng Kolmogorov–Arnold). Kiến trúc này cho phép tổng hợp và xử lý đồng thời ba nguồn dữ liệu hình ảnh, văn bản và âm thanh trong một không gian biểu diễn thống nhất. Thực nghiệm trên tập dữ liệu ICH tự thu thập cho thấy phương pháp đề xuất vượt trội so với các mô hình baseline, trong đó sự kết hợp DenseNet121 + FastKAN đạt độ chính xác cao nhất 94.35\% và các mô hình gọn nhẹ như MobileNetV3 Small hay ShuffleNetV2 1.0x ghi nhận mức cải thiện trên 8\%. Ngoài ra, việc tích hợp FastKAN giúp mô hình xấp xỉ tốt hơn các biên quyết định phi tuyến phức tạp so với MLP truyền thống. Tính thực tiễn của nghiên cứu còn được khẳng định thông qua việc xây dựng và triển khai hệ thống web ICH Predictor, hỗ trợ tra cứu và trực quan hóa quy trình phân loại dữ liệu ICH một cách thuận tiện và hiệu quả.\vspace{0.3cm}

\textbf{Từ khóa:} Di sản văn hóa phi vật thể, Phân loại đa phương thức, Mạng noron đồ thị, MGAT, FastKAN, Học sâu.